\newpage
% \color{white}
% \pagecolor{black}

\ifdefined\useChapters
\chapter{A busca}
\else
\chapter{}
\fi

O chão de terra ainda guarda o calor do dia, e o ar noturno sopra úmido entre as árvores. A fogueira improvisada lança sombras longas nos troncos. Sento-me ao lado dela com o caderno aberto sobre os joelhos — o mesmo em que escrevo desde que tudo começou. O papel já está amassado, sujo de terra e fuligem, mas ainda é o único lugar onde posso tentar organizar o que vi.

A discussão sobre viagens no tempo era divertida quando cabia nos livros. Sublinhei teorias sobre paradoxos e contínuo temporal como quem avalia regras de um jogo que jamais precisará jogar. Agora encosto a ponta do lápis no papel e lembro de conversar comigo mesmo no mercadinho. A mina quebra. Eu estive lá; qualquer teoria que eu escrever pode ser também a lembrança de algo que já provoquei.

Escrevo devagar, tentando dar forma ao pensamento:  
“Mudar o futuro é sempre uma incógnita. Se temos livre-arbítrio, o futuro não deveria estar escrito.”  
A frase fica torta na página. Sopro o carvão do lápis e continuo.  

Escrevo uma segunda frase: “Se eu vi o futuro, ele já existe?” Risco. Escrevo: “Se existe, por que tentar mudar?” Risco também. Gosto de pensar que o futuro continua sendo desenhado pelas nossas ações, mas a realidade não se importa com a versão que me tranquiliza.

Depois de ver aquele futuro apocalíptico — pessoas dominadas pelo poder das próprias habilidades, destruindo sem pensar — confesso que minha crença no bem ficou abalada. Sempre imaginei que alguns se perderiam, que o poder corromperia um ou outro, mas o que vi foi um mundo inteiro queimando de raiva e prepotência.  

O fogo ainda parece preso dentro de mim, como se tivesse respirado cinza demais.  
Fecho o caderno por um instante e encaro a fogueira. A chama se retorce com o vento, e por um segundo parece responder aos meus pensamentos. 

Enquanto estive naquele futuro, vi habilidades incríveis.  
Um rapaz atravessava paredes, mas não como um fantasma — o corpo dele mudava, se tornava parte da matéria que tocava.  
Outro manipulava bolas de energia como se brincasse com luz.  
Uma menina se movia de um ponto a outro num piscar de olhos — teletransporte, se é que essa palavra explica.  
E um rapaz flutuava acima de tudo, sereno, observando.  
De todas, a dele parecia a mais inofensiva. Talvez porque eu o reconhecesse — o mesmo que me observava da floresta quando fui perseguido.  
Mesmo ali, entre o fogo e a loucura, ele só olhava.  
Isso me marcou.  

Agora, pensando melhor, o poder mais perigoso talvez seja o de apenas observar e chamar a própria distância de neutralidade.

Respiro fundo e volto a escrever.  
“Depois de tanto tempo desdobrado, aprendi que, se eu me concentrar, posso ir a qualquer lugar que conheça — qualquer ponto do espaço e do tempo. Mas não consigo tocar o mundo. Sou só espírito.”  

As palavras parecem frias demais para o que sinto. Mesmo assim, desenho ao lado delas meu corpo junto à fogueira e uma segunda figura acima das árvores. Estar em todos os lugares e em nenhum parece bonito no papel. Na prática, significa deixar o corpo indefeso enquanto uma versão de mim decide onde entrar sem ser convidada.

Penso nas coisas que consegui fazer enquanto estava fora do corpo — criar ilusões, viajar no tempo, volitar.  
Talvez essas sejam outras habilidades, apenas mais acessíveis quando não temos o corpo pra nos limitar.  
E se for isso?  
E se o corpo for o que nos prende?  
Um véu que cobre o que realmente somos?

A fogueira estala. Um graveto explode em faíscas que sobem e desaparecem.  
Olho pra elas e penso que talvez seja isso: somos faíscas contidas, acreditando que somos o tronco inteiro.

Se estou certo, então todos podem.  
Todos nós temos o mesmo potencial.  
E o que nos impede é só a crença no impossível.

Fecho o caderno e me deito, usando a mochila como travesseiro.  
O céu, visto entre as copas, está cheio de estrelas.  
Penso se, quando dormimos, não nos desdobramos também — se os sonhos não são, na verdade, viagens por outros tempos, outras versões de nós mesmos.  

O corpo cede, pesado.  
Antes de adormecer, repito pra mim mesmo: amanhã começo minha busca.  
Vou atrás do volitador.  
Quero entender como ele faz.  
Talvez, quando descobrir, eu também consiga.  
E se conseguir...  
talvez descubra um outro universo escondido dentro deste.

O vento apaga a fogueira aos poucos, e o silêncio da floresta cobre tudo.  
Só o som do meu coração me lembra que ainda estou aqui — por enquanto.
